% sec-05: outline item E.  Figures 13 (mermaid), 14 (graphviz), 15 (plantuml).
% Table 13 is the single artifact most likely to convert an inquiry into a call.
\section{Twelve Milestones a Program Officer Can Audit}

\subsection{What an Auditable Milestone Is}

\draftinstr{Four sentences defining the contract this section offers. A
milestone is auditable only if it carries four things: an evidence artifact that
exists as a file, a cost that sums into its phase, a stop condition stated
before the work begins, and a named funder mechanism that pays for it. State
that all twelve rows carry all four, and that the two phase totals are exact.}

\figslot{mermaid-type / gantt}{7.4}{Thirty-three months, twelve milestones, and
the artifact date on\\each. Phase I holds five at 306,000 dollars; Phase II
holds seven\\at 1,300,000. The gate at month nine is the only hard boundary.}

\subsection{The Twelve Rows}

\draftinstr{Set Table 13, the central artifact of the paper. Six columns: ID,
months, milestone, cost, evidence artifact, stop condition. Twelve rows. The
first five must sum to \$306,000 and the last seven to \$1,300,000. Every row's
data is in
\appfile{funding/capitalization-plan/mermaid/fig-13-twelve-milestone-calendar.md}
and every stop condition must be a number or a signed instrument, never a
judgement. Because this table is wide, cut it as
\texttt{p\{0.8cm\} p\{1.4cm\} X p\{1.6cm\}} for the first four columns and set
the evidence artifact and stop condition as a second table, Table 14, keyed by
the same twelve IDs, rather than compressing six columns into one measure.}

\draftinstr{Set Table 14: twelve rows keyed by milestone ID, three columns for
evidence artifact, stop condition, and the funder mechanism that pays. The
mechanism column takes four values only: SBIR Phase I, SBIR Phase II, the
\$2,104,000 delta, and contributed. Source:
\appfile{funding/capitalization-plan/graphviz/fig-09-work-package-dag.gv.md}
for the mechanism assignment, which must agree with WP1 to WP12.}

\subsection{How a Milestone Is Closed}

\figslot{plantuml-type / activity with fork}{10.4}{One milestone, three
concurrent evidence branches, and the join\\that has to close before a program
officer can audit anything.\\The replay is the only step that can send a
milestone backwards.}

\draftinstr{After Figure 15, one paragraph on why all three branches must join
before an audit can begin: the hash proves the artifact has not moved, the
source record proves it describes something that happened, and the monitoring
report proves a third party looked at both. Any one alone is a claim, not
evidence. Source:
\appfile{funding/capitalization-plan/plantuml/fig-15-milestone-evidence-activity.puml.md}
and \appfile{funding/pdac-funding-applications/final-apply/sections/sec-06-physical-ai-governance.tex}.
Cite \texttt{ecfr2026part312records}.}

\subsection{How the Programme Stops}

\figslot{graphviz-type / fault tree}{8.6}{One top event, five halt points, and
ten basic events below.\\Three are AND gates, so no single failure stops the
programme.\\The two OR branches are the month-nine gate and ownership test.}

\draftinstr{After Figure 14, two paragraphs. The first names the two OR branches
as the honest part of the figure: H1 and H5 are the only single points of
failure and neither is technical. The second explains why the top event is a
stop \emph{without a deposited archive} rather than a stop, and walks the
survival band: at every halt point the archive is already deposited for every
closed milestone, so reaching the top event requires a halt plus a failure to
deposit. Forward-reference Figure 20. Source:
\appfile{funding/capitalization-plan/graphviz/fig-14-stop-condition-fault-tree.gv.md}.}
